\section{Basic Results on Random Vectors \label{appendix:basic}}

The following is a well-known estimate for the tail of a Gaussian.

\begin{proposition}
When $Z$ is a unit Gaussian, we have
$$
\ln \Pp(Z \ge a) = - \frac 1 2 a^2 - \ln a - \frac 1 2 \ln 2 \pi + o(1).
$$
for $a \to + \infty.$
\end{proposition}

In fact, the leading-order term $- \frac 1 2 a^2$ is an upper bound on $\ln \Pp(Z \ge a).$ One simple way to establish such estimates on tail probabilities is via the Chernoff inequality: for any random variable $X$ and $\lambda > 0,$
$$
\Pp(X \ge a) \le e^{-\lambda a} \E e^{\lambda X} = \exp(- \lambda a + K_X(\lambda))
$$
where $K_X(\lambda) = \ln \E e^{\lambda X}$ is the cumulant generating function. In general, putting $\lambda = a$ in the Chernoff inequality for a random variable $X$ with $K_X(\lambda) \le \lambda^2 / 2$ gives
$$
\Pp(X \ge a) \le \exp(-a^2/2),
$$
which is called a sub-Gaussian tail bound.

Let $X_1, \dots, X_k$ be independent Rademacher random variables, each uniformly distributed on $\{-1, 1\}.$ Then $K_{X_i}(\lambda) = \ln \cosh \lambda \le \lambda^2 / 2,$ so
$$
K_{R_k}(\lambda) = \sum_{i=1}^k K_{X_i}(\lambda) \le k \lambda^2 / 2
$$
where $R_k = X_1 + \dots + X_k.$ This proves the following.

\begin{proposition}[Chernoff bound for Rademacher sums]
\label{prop:rademacher-chernoff}
For $R_k$ a sum of $k$ independent Rademacher variables and $t > 0,$
$$
\Pp(R_k \ge t) \le \exp\left(-\frac{t^2}{2k}\right).
$$
\end{proposition}

This is enough to prove \Cref{prop:many-ortho}, which we restate here for convenience.

\begin{proposition*}
	Let $d > 2 \epsilon^{-2} (2 \ln N + \ln p^{-1}),$ and let
	$$
		\{ F_1, \dots, F_N \} \subseteq \{ - 1/\sqrt d,  1 / \sqrt d \}^d
	$$
	be random vectors with independent, uniformly distributed entries. Then $\lvert \langle F_i, F_j \rangle \rvert < \epsilon$ for all $i \neq j$ with probability at least $(1 - p).$
\end{proposition*}

\begin{proof}
    Each inner product $\langle F_i, F_j \rangle$ is distributed like $R_d / d$ where $R_d$ is a sum of $d$ Rademacher variables. By the Chernoff bound, $\Pp(\langle F_i, F_j \rangle \ge \epsilon) = \Pp(R_d \ge d\epsilon) \le \exp(- d \epsilon^2 / 2).$ By symmetry and a union bound over all $\binom N 2 < N^2 / 2$ pairs,
    $$
    \Pp(\exists\, i \neq j : \lvert \langle F_i, F_j \rangle \rvert \ge \epsilon) \le N^2 \exp(- d \epsilon^2 / 2).
    $$
    This is at most $p$ when $d \ge 2 \epsilon^{-2} (2 \ln N + \ln p^{-1}).$
\end{proof}

The interested reader should also compare this result to the Johnson-Lindenstrauss lemma, which is proved in a very similar way. (See \cite{dasgupta_elementary_2003} for a proof, or the last section of \cite{foucart_sparse_2013} for a discussion of the JL lemma with some broader context.)

Now, consider the inner product of two $d$-dimensional codewords $X_i, X_j$ drawn from a random spherical dictionary. (That is, $X_i$ and $X_j$ are uniformly distributed on the unit sphere in $d$ dimensions.) Write $b(d, t)$ for the probability that $\langle X_i, X_j \rangle \ge t.$ This number can also be viewed as fraction of the unit sphere occupied by the spherical cap
$$
\{ (v_1, \dots, v_d) : v_1^2 + \dots + v_d^2 = 1, \; v_1 \ge t \}
$$
Using the fact that normalizing a Gaussian random vector gives a uniform draw from the unit sphere, it is possible to establish that
$$
b(d, t) \le (1 + o(1)) \exp\left(- \frac {t^2 d} 2 \right)
$$
using Chernoff inequalities. However, for the proof of \Cref{prop:necessary-cond}, we also need a \textit{lower} bound on this tail probability. This is accomplished by the following result, which is a corollary of some ideas from Lecture 2 of \cite{levy_elementary_1997}.
\begin{proposition} \label{prop:ball}
For all $t \in (0, 1),$
$$
\exp\left(-\frac{t^2 d}{2(1-t^2)}\right) \le b(d, t) \le \exp\left(-\frac{t^2 d}{2}\right)
$$
\end{proposition}
